\section{Scope, conventions, and the proof ledger}
\label{scope:sec}

The source corpus already contains substantial all-order mathematics.
In particular, the incoming reports \cite{SourceCanonical,SourceHigher} establish Gamma-quotient harmonic
generators, arbitrary digamma-power algorithms, the fifth moment,
continuous-depth Lerch constructions, and several Stieltjes collision
formulas.  Repeating those results as new would obscure the actual
progress.  We use them as a baseline and concentrate on four concrete
tasks: resolving a recorded conjecture on an infinite family, proving a
recorded numerical polylogarithm relation, finding closed coefficient
formulas inside the higher Gamma hierarchy, and completing explicit
higher-moment and spectral-derivative calculations.

\subsection{What is proved here, and what is inherited}

\begin{center}
\small
\begin{tabularx}{\textwidth}{@{}p{.26\textwidth}X@{}}
\toprule
Topic & Precise status in this continuation \\
\midrule
Gaussian Family--2 & The numerical source equation is proved exactly,
and reduced to $\zeta(3)$, $\pi\beta(2)$, and $\pi^2\log2$.
The general parity theorem is established prior work. \\
Adjacent even tails & The first-$q$ conjecture is proved for all
$p=q-2\ge2$.  The exact determinant and reconstruction identities
are new developments relative to the inspected reports. \\
Differential normal form & The functional-monomial decomposition is
classical.  Its Gamma coefficient sectors and explicit endpoint
specialization are developed and proved here. \\
Quadratic Hurwitz products & Explicit integer finite parts and every
order-derivative jet follow from the classical product kernel, with
the local coordinate correction proved. \\
Sixth digamma moment & A complete explicit formula is supplied in
height-one Laurent slopes and one specified, convergent depth-at-most-three
zeta-derivative combination.  Its further arithmetic reduction remains open. \\
$S_6$, revised $S_8$, larger gaps & These remain conjectural or open.
Neither the parity argument nor finite numerical evidence resolves them. \\
\bottomrule
\end{tabularx}
\end{center}

The distinction between an \emph{explicit coordinate formula} and a
\emph{reduction to a prescribed constant alphabet} matters.  Defining a
constant by an unexplained divergent sum would not complete a moment
calculation.  Our remaining mixed coordinate is instead given by an
absolutely convergent logarithmic series, a finite multiple-zeta
expression, an exact value at spectral order zero, and independently
checked analytic formulas.  We do not assert that it is irreducible or
independent of the other constants.

\subsection{Harmonic, polylogarithmic, and Laurent conventions}

For $n\ge0$ and $r\ge1$,
\[
 H_n^{(r)}=\sum_{k=1}^n k^{-r},\qquad H_n=H_n^{(1)},
 \qquad H_0^{(r)}=0.
\]
For a word $\boldsymbol a=(a_1,\ldots,a_d)$ of positive integers,
\[
 H_n(\boldsymbol a)=\sum_{n\ge k_1>\cdots>k_d\ge1}
             \frac1{k_1^{a_1}\cdots k_d^{a_d}},
 \qquad H_n(\varnothing)=1.
\]
All multiple zeta functions and multiple polylogarithms in this article
use \emph{decreasing} indices:
\[
 \zeta(s,\boldsymbol a)=\sum_{n\ge1}\frac{H_{n-1}(\boldsymbol a)}{n^s},
 \qquad
 \Li_{p,q}(x,y)=\sum_{n>m\ge1}\frac{x^ny^m}{n^pm^q}.
\]
Outside their initial domains, multiple zeta functions mean their
meromorphic continuations.  Where a displayed finite combination is
holomorphic although individual terms have poles, the combination is
formed \emph{before} its Taylor coefficient is taken.

We use the conventional generalized Stieltjes constants
\begin{equation}\label{scope:stieltjes}
 \zeta(1+u,a)=\frac1u+
       \sum_{n\ge0}\frac{(-1)^n\gamma_n(a)}{n!}u^n,
 \qquad \gamma_0(a)=-\psi(a),\quad \gamma_n=\gamma_n(1).
\end{equation}
Thus the coefficient of $u$ at $a=1$ is $-\gamma_1$, not
$\gamma_1$.  A second family used below is
\begin{equation}\label{scope:lambda}
 \lambda_d=[u^1]\zeta(1+u,\{1\}^{d-1}),\qquad d\ge1.
\end{equation}
Here $[u^1]$ means the Laurent coefficient; the function may have a pole
of order $d$.  In particular $\lambda_1=-\gamma_1$.
The higher $\lambda_d$ are strict-depth Laurent slopes, not a relabeling
of the ordinary $\gamma_d$.  The local notation $\lambda_2$ for a
dilogarithm imaginary part in the source Family--2 comparison is
explicitly confined to that comparison.

The basic differential identities are
\begin{equation}\label{scope:hurwitz-differentiation}
 \partial_a\zeta(s,a)=-s\zeta(s+1,a),\qquad
 \psi^{(r)}(a)=(-1)^{r+1}r!\zeta(r+1,a)\quad(r\ge1).
\end{equation}
They follow directly from the defining series in its absolute
convergence domain and then by continuation; see also DLMF~25.11
\cite{DLMF}.  They fix the signs when translating the harmonic-tail
identities into polygamma products.

For later comparison with antiderivatives, coefficient extraction in
\eqref{scope:hurwitz-differentiation} gives the classical formulas
\begin{align}
 \gamma_n'(a)
  &=(-1)^{n+1}
     \left\{\left.\partial_s^n\zeta(s,a)\right|_{s=2}
      +n\left.\partial_s^{n-1}\zeta(s,a)\right|_{s=2}\right\},
       \label{scope:stieltjes-derivative}\\
 \int_{a_0}^{a}\gamma_n(v)\dd v
  &=\frac{(-1)^{n+1}}{n+1}
    \left.\partial_s^{n+1}
       \{\zeta(s,a)-\zeta(s,a_0)\}\right|_{s=0}.
       \label{scope:stieltjes-primitive}
\end{align}
The second term in braces in \eqref{scope:stieltjes-derivative} is
absent for $n=0$.  The derivative formula comes from
$\partial_a\zeta(1+u,a)=-(1+u)\zeta(2+u,a)$.
For the primitive use
$\partial_a\zeta(u,a)=-u\zeta(1+u,a)$, then differentiate
$n+1$ times at zero.  These formulas are included to connect the new
spectral-jet identities with the conventional Stieltjes normalization;
no priority claim is made for them.

\subsection{The endpoint coordinate is part of the definition}

If a function on $(0,1]$ has a finite principal Laurent part at zero,
we define
\begin{equation}\label{scope:finitepart}
 \FP\int_0^1 f(x)\dd x
   =\CT_{\epsilon=0}\int_\epsilon^1 f(x)\dd x.
\end{equation}
The same notation applies to finite sums of powers times logarithms,
followed by an integrable remainder.  The constant coefficient is taken
after all negative powers and all divergent logarithms are removed.
The cutoff is always the displayed coordinate $x=\epsilon$.
In particular,
\[
 \FP\int_0^1x^{-j}\dd x=-\frac1{j-1}\ (j\ge2),\qquad
 \FP\int_0^1x^{-1}(\log x)^a\dd x=0\ (a\ge0).
\]
For $h$ meromorphic near $[0,1]$ with no poles on $(0,1]$,
\begin{equation}\label{scope:boundary}
 \FP\int_0^1h'(x)\dd x=h(1)-[x^0]h(x).
\end{equation}
This is an exact coefficient identity obtained by integrating from
$\epsilon$ to $1$.  It does not replace $[x^0]h$ by the generally
nonexistent value $h(0)$.

Finally, throughout the tail-rank and differential-polynomial sections,
``rank'' and ``dimension'' refer to coefficient spaces over $\Q$.
For example, using distinct formal symbols for $\zeta(0),\zeta(2),\ldots$
does not assert their algebraic independence.  Likewise the image of a
derivation is a vector subspace, generally not an ideal.
